\section{Materiales y métodos}

\hspace{1cm}\textbf{Extracción de las secuencias en del grupo K01220 en KEGG}

Se utilizará la API de KEGG \parencite{Kanehisa_2004, Kanehisa_2022} para la extracción automática de datos. La documentación oficial está disponible en \url{https://www.kegg.jp/kegg/rest/}. Se desarrollarán scripts en Python (versión 3.12.7), implementados en notebooks de Jupyter (.ipynb), con el objetivo de:
\begin{itemize}
	\item Extraer las secuencias del grupo K01220.
	\item Realizar la limpieza de los datos.
	\item Generar archivos de salida en formato FASTA (.fas).
	\item Registrar el procesamiento y análisis de los datos.
\end{itemize}
Todo este código estará disponible en la liga \url{https://github.com/Electrocyte96/tesina/tree/main/scripts}

\hspace{1cm}\textbf{Alineamiento múltiple de secuencias}
Una vez obtenidas las secuencias en formato FASTA se utilizará el software MEGA11 \parencite{Tamura_2021} para realizar un alineamiento múltiple con el algoritmo MUSCLE \parencite{Edgar_2004} 

\hspace{1cm}\textbf{Búsqueda de dominios conservados con dbCAN}

Sobre las secuencias sin alineamiento previo se usará la herramienta dbCAN \parencite{Zheng_2023, Zhang_2018, Huang_2017}. Esta plataforma permite identificar clústeres de genes relacionados con carbohidratasas (CGCs), también realizar predicción de substratos según una determinada secuencia de aminoácidos y  permite la anotación de carbohidratasas utilizando perfiles de dominio derivados de la base de datos CAZy (Carbohydrate-Active enZYmes Database). La anotación se llevó a cabo utilizando tres estrategias: búsquedas basadas en modelos ocultos de Markov mediante pyHMMER \parencite{Larralde_2023}, el cual es una implementación en Python/Cython de HMMER \parencite{Eddy_2011}; alineamientos de secuencia mediante DIAMOND \parencite{Buchfink_2021}; y clasificación en subfamilias mediante dbCAN-sub. Estas estrategias permiten detectar dominios conservados y asignar familias CAZy a cada secuencia analizada.

Como resultado de este proceso podremos obtener una serie de archivos en formato de Valores Separados por Tabulaciones (.tsv) los cuales describirán los motivos reconocidos por esta herramienta. 

\hspace{1cm}\textbf{Generación de árboles de similitud}

\parencite{Minh_2020, Wong_2025}

\hspace{1cm}\textbf{Análisis estructural con DALI}

Sobre las secuencias de interés se usará la herramienta en línea DALI server \parencite{Holm_2023} para el análisis estructural de proteínas basándose en un alineamiento de matrices de distancia \parencite{Holm_2019}. Es utilizada ampliamente por cristalógrafos para buscar parentescos de nuevas estructuras proteicas con otras proteínas previamente conocidas y almacenadas en las bases de datos del Protein Data Bank (PDB) y de AlphaFold Database \parencite{Holm_2022} a su vez de apoyarse en Pfam \parencite{Mistry_2020} para la anotación de proteínas.   